% Section 4: The six-state carrier and the derived Clifford relation
\label{sec:six-state-body}

Three readings, each two-dimensional over the reals, give six real dimensions.  The
temptation is to equip each pair with a complex structure and declare the result
$\C^{3}$.  This section shows that the correct carrier is not a sum of three copies of
anything: it is the six-dimensional space
\[
  W=\Lambda^{1}E\oplus\Lambda^{2}E,\qquad \dim_{\R}E=3,
\]
in which each reading is paired with the \emph{oriented relation of the other two},
and in which the complex unit appears exactly once.  The section then shows that on a
declared finite cut interval this structure is not postulated but generated.

\subsection{The carrier and its paired planes}

\begin{definition}[Six-state carrier]\label{def:six-state-carrier}
Let $E$ be a three-dimensional real vector space with basis $T,S,R$ representing the
temporal, spatial, and relational readings.  Put
\[
  W=\Lambda^{1}E\oplus\Lambda^{2}E=\langle T,S,R\rangle\oplus\langle SR,RT,TS\rangle ,
\]
so $\dim_{\R}W=6$.  The three \emph{aspect--opposite-relation planes} are
\[
  W_T=\langle T,SR\rangle,\qquad W_S=\langle S,RT\rangle,\qquad W_R=\langle R,TS\rangle .
\]
\end{definition}

The pairing is the content of the definition: a reading is not paired with a second
copy of itself but with the oriented plane spanned by the other two.  This is forced
by the dimension count once one asks that each reading be recoverable from its own
opposite relation, and it has an immediate consequence for the complex structure.

\begin{proposition}[One complex unit, not three]\label{prop:one-complex-unit}
Suppose the generators satisfy
\[
  T^{2}=S^{2}=R^{2}=1,\qquad \varepsilon_i\varepsilon_j+\varepsilon_j\varepsilon_i=0
  \ \ (i\neq j),
\]
and set $\Omega=TSR$.  Then $\Omega$ is central and
\[
  \Omega^{2}=-1 .
\]
Consequently $J=\Omega\cdot$ preserves $W$, swaps $\Lambda^{1}E$ with
$\Lambda^{2}E$ by cyclic pairing
\[
  J(T)=SR,\quad J(SR)=-T;\qquad J(S)=RT,\quad J(RT)=-S;\qquad
  J(R)=TS,\quad J(TS)=-R,
\]
and on each plane $W_T,W_S,W_R$ restricts to
$\begin{pmatrix}0&-1\\ 1&0\end{pmatrix}$.  Hence
$\chi_J(\lambda)=(\lambda^{2}+1)^{3}$.
\end{proposition}

\begin{proof}
\emph{Centrality.}  A generator $T$ commutes with $SR$ up to $(-1)^{2}=+1$, so
$\Omega$ lies in the centre.

\emph{The square.}  Moving the fourth factor of $\Omega^{2}=T\,S\,R\,T\,S\,R$ to the
front costs the two transpositions past $R$ and $S$, hence $(-1)^{2}=+1$:
\[
  T\,S\,R\,T\,S\,R = T^{2}\,S\,R\,S\,R = S\,R\,S\,R ,
\]
using $T^{2}=1$.  Moving the third factor past the fourth costs one transposition:
\[
  S\,R\,S\,R = -\,S^{2}R^{2} = -1 .
\]
Hence $\Omega^{2}=-1$.

\emph{Pairing.}  For the first plane, $RT=-TR$ gives
\[
  \Omega T=T\,S\,R\,T=T\,S\,(-T\,R)=-(T\,S\,T)R=(T\,T\,S)R=T^{2}SR=SR,
\]
and then $\Omega(SR)=T\,S\,R\,S\,R=-T\,S^{2}R^{2}=-T$.  The other two planes follow by
cyclic symmetry.  On the plane spanned by $T$ and $SR$ the action is therefore
$T\mapsto SR\mapsto -T$, which is the displayed matrix.

\emph{Characteristic polynomial.}  $J$ acts as the same $2\times2$ rotation on each of
the three planes, so
$\chi_J(\lambda)=\prod_{i}\det\bigl(\lambda I-J|_{W_i}\bigr)=(\lambda^{2}+1)^{3}$.
\end{proof}

\begin{calibration}[The algebraic form is an assumption]\label{cal:algebraic-form}
Proposition~\ref{prop:one-complex-unit} takes the anticommutation relations as
\emph{hypotheses}.  It is a bounded algebraic calibration: it shows what the carrier
does once such generators exist.  \S\ref{sec:six-state-derivation} supplies the
derivation.
\end{calibration}

\begin{proposition}[Chirality sign law]\label{prop:sign-law}
For every permutation $\pi$ of the three directions,
\[
  P_\pi\,\Omega\,P_\pi^{-1}=\sgn(\pi)\,\Omega .
\]
In particular exchanging temporal and spatial directions gives
$\Omega\mapsto-\Omega=\Omega^{-1}$, so cyclic permutations preserve the orientation
expression and transpositions reverse it.
\end{proposition}

\begin{proof}
$P_\pi$ permutes the generators, hence multiplies the word $TSR$ by the sign of the
permutation.  Invariantly: a transposition reverses the induced orientation of $E$,
and $\Omega$ is the orientation class of $E$ followed by the volume of
$\Lambda^{2}E$ under $\Omega$'s own duality.
\end{proof}

\begin{remark}
Proposition~\ref{prop:sign-law} is the carrier-level form of
Definition~\ref{def:chirality}: the central sign is exactly what a transposition
produces, and it is invisible to any projective reading.
\end{remark}

\subsection{The signature trichotomy}

If the three directions are granted arbitrary signs, the complex unit becomes
conditional.

\begin{proposition}[Signature trichotomy]\label{prop:signature-trichotomy}
With $T^{2}=\varepsilon_T$, $S^{2}=\varepsilon_S$, $R^{2}=\varepsilon_R$ and
$\varepsilon_i\in\{-1,0,+1\}$,
\[
  \Omega^{2}=-\varepsilon_T\varepsilon_S\varepsilon_R ,
\]
so the carrier is elliptic when the product is positive, hyperbolic when negative, and
parabolic when the product vanishes.
\end{proposition}

\begin{proof}
Each $T^{2}$ contributes a factor $\varepsilon_T$ to the reduction performed in the
proof of Proposition~\ref{prop:one-complex-unit}; the signs multiply.
\end{proof}

\begin{warning}[Temporality is not a metric sign]\label{warn:temporality}
A Lorentzian reading with exactly one negative square forces the split case.  Hence
semantic temporality must \emph{not} be identified with a negative metric direction;
doing so silently converts the trichotomy of \S\ref{sec:scalar-trichotomy} into an
unannounced choice of $\kappa=1$.
\end{warning}

\subsection{Deriving the generators from cuts}\label{sec:six-state-derivation}

The material above is algebraic.  The geometric content of this section is that on a
declared finite cut interval the two operations needed are already present, so the
Clifford relation is \emph{derived}.

\begin{construction}[Local cut interval]\label{con:cut-interval}
Let $U$ be a finite set of events carrying the completion order of a checked acyclic
program, and let $g_1,g_2,g_3$ be three pairwise independent enabled events
generating a local Boolean cut interval.  Two operations are available on cut data:
\begin{enumerate}
\item \emph{oriented interchange}, which exchanges two independent enabled events and
  records the sign of the exchange, giving endomorphisms $\varepsilon_i$ on the cut
  cochains;
\item \emph{cut contraction}, which reads whether an event belongs to a cut, giving
  $\iota_i$ determined by the membership pairing
  \[
    dx_i(g_j)=\delta_{ij}.
  \]
\end{enumerate}
\end{construction}

\begin{theorem}[Endogenous Clifford relation]\label{thm:endogenous-clifford}
On the local cut interval of Construction~\ref{con:cut-interval} the following hold.
\begin{enumerate}
\item Oriented interchange yields the exterior relations
  $\varepsilon_i\varepsilon_j+\varepsilon_j\varepsilon_i=0$ for $i\neq j$.
\item With membership pairing $dx_i(g_j)=\delta_{ij}$, the operators
  $c_i=\varepsilon_i+\iota_i$ satisfy
  \[
    c_ic_j+c_jc_i=2\delta_{ij}I .
  \]
\item Hence $\Omega=c_1c_2c_3$ satisfies $\Omega^{2}=-I$, and on
  $W=\Lambda^{1}E\oplus\Lambda^{2}E$ the orientation action has
  $\chi_J(\lambda)=(\lambda^{2}+1)^{3}$ with three aspect--opposite-face planes.
\end{enumerate}
\end{theorem}

\begin{proof}[Proof sketch and status]
Items (i) and (ii) are the Chevalley construction of a Clifford algebra from an
exterior algebra together with a contraction satisfying the duality
$dx_i(g_j)=\delta_{ij}$; what the construction adds is that both operations are
\emph{supplied by the cut data} rather than assumed.  Item (iii) is then
Proposition~\ref{prop:one-complex-unit} applied to the derived generators.
The derivation is checked exhaustively on one local three-cube and is reported as a
bounded endogenous derivation, not as a theorem about arbitrary cut intervals.
\end{proof}

\begin{theorem}[Interchange alone is insufficient]\label{thm:interchange-insufficient}
Oriented interchange by itself yields only a nilpotent exterior volume: it produces
$\varepsilon_1\varepsilon_2\varepsilon_3$ satisfying
$(\varepsilon_1\varepsilon_2\varepsilon_3)^{2}=0$, and no complex structure.  The
contraction $\iota_i$ is indispensable.
\end{theorem}

\begin{proof}
Wedge products of degree three in a three-dimensional exterior algebra have square
zero because any repeated factor annihilates the product.  Without a contraction
there is no degree-lowering operation, so the Chevalley combination $c_i=\varepsilon_i
+\iota_i$ cannot be formed.
\end{proof}

\begin{remark}[Why this matters]
Theorems~\ref{thm:endogenous-clifford} and~\ref{thm:interchange-insufficient}
together answer the natural objection to \S\ref{sec:scalar-trichotomy}: the complex
structure is not smuggled in with the carrier.  It is produced by two operations that
the cut interpretation already needs, and the second of them is doing essential work.
\end{remark}

\subsection{Filtrations from a designated causal line}

One further structure is available, and it is important not to over-read it.

\begin{proposition}[Six-state filtration]\label{prop:filtration}
Designating one causal line refines the six-state carrier as
\[
  6=1+2+2+1 .
\]
The two readings $(1,2,3)$ and $(1,3,2)$ obtained by regrouping the transverse pairs
differ by an orientation-unit transfer and are \emph{not} a chiral pair.
\end{proposition}

\begin{proof}[Status]
The count is elementary from
Definition~\ref{def:six-state-carrier}.  The absence of chirality is the content: see
Theorem~\ref{thm:regrouping-not-chirality}.
\end{proof}

\begin{theorem}[No-go: regrouping is not chirality reversal]\label{thm:regrouping-not-chirality}
Orientation-unit transfer changes block membership but not $\Omega$, whereas an actual
odd exchange of two transverse directions satisfies $P\Omega P^{-1}=-\Omega$.  Hence
the passage $(1,2,3)\leftrightarrow(1,3,2)$ is a change of filtered reading and not a
chiral pair.  Consequently an ordinary scalar spectrum cannot distinguish the two
readings, while the filtration can.
\end{theorem}

\begin{proof}
The two facts are verified separately: the transfer preserves the orientation word by
Proposition~\ref{prop:sign-law} applied to an even permutation, while an odd exchange
is a transposition and therefore reverses it.  A scalar spectrum depends only on
$\Omega^{2}$ and hence cannot see the sign at all.
\end{proof}

\begin{remark}
Theorem~\ref{thm:regrouping-not-chirality} is a small result with a large moral, which
recurs in \S\ref{sec:residual-catalogue}: quantities built from squares, norms, or
spectra are blind precisely to the central sign that
Definition~\ref{def:chirality} isolates.
\end{remark}
